\documentclass[11pt]{amsart}

\usepackage[T1]{fontenc}
\usepackage{lmodern}
\usepackage{microtype}
\usepackage{amsmath,amssymb,amsthm}
\usepackage{booktabs}
\usepackage[hidelinks]{hyperref}

\hypersetup{
  pdftitle={A Five-Vertex Counterexample to Tang's Bunkbed Self-Avoiding-Walk Question},
  pdfsubject={Self-avoiding walks on bunkbed graphs}
}

\newtheorem{theorem}{Theorem}
\newtheorem{proposition}[theorem]{Proposition}
\theoremstyle{remark}
\newtheorem{remark}[theorem]{Remark}

\newcommand{\SAW}{\mathcal S}

\title[A five-vertex bunkbed counterexample]
{A Five-Vertex Counterexample to Tang's Bunkbed Self-Avoiding-Walk Question}
\date{}

\begin{document}

\begin{abstract}
Tang asked whether, for every finite connected simple graph $G$, the
number of self-avoiding walks in $G\square K_2$ from $u_0$ to $v_0$
is at most the number from $u_0$ to $v_1$ whenever $uv$ is not a cut
edge. We answer this question negatively. For the five-vertex house
graph and a cycle edge $ad$, exact enumeration gives
\[
  |\SAW(a_0,d_0)|=60>59=|\SAW(a_0,d_1)|.
\]
The base graph is $2$-connected, and five is the minimum possible
order among counterexamples with an adjacent non-cut edge. We also
give a $3$-connected counterexample on seven vertices.
\end{abstract}

\maketitle

\section{A minimum-order counterexample}

For a finite connected simple graph $G$, write $x_i=(x,i)$ for a
vertex of the bunkbed graph $G\square K_2$. For vertices $p,q$ of the
bunkbed graph, let $\SAW(p,q)$ denote the set of self-avoiding walks
from $p$ to $q$. Tang asked whether
\begin{equation}
  |\SAW(u_0,v_0)|\le |\SAW(u_0,v_1)|
  \label{eq:tang}
\end{equation}
holds whenever $(u,v)$ is not a cut edge of $G$
\cite[Question~3.3]{Tang}. The example below uses an adjacent pair,
so it gives a negative answer independently of any convention
concerning nonadjacent pairs.

\begin{theorem}
Let
\[
  V(G)=\{a,b,c,d,e\},\qquad
  E(G)=\{ab,ad,ae,bc,be,cd\}.
\]
Then $G$ is $2$-connected, $ad$ is not a cut edge, and
\[
  |\SAW(a_0,d_0)|=60>59=|\SAW(a_0,d_1)|.
\]
Consequently, \eqref{eq:tang} is false. Moreover, no counterexample
with an adjacent non-cut edge has fewer than five vertices.
\end{theorem}

\begin{proof}
The graph $G$ is the house graph: it is the union of the cycle
$a b c d a$ and the triangle $a b e a$. Deleting any vertex leaves a
connected graph, so $G$ is $2$-connected. In particular, the cycle
edge $ad$ is not a cut edge.

We use an exact last-step recurrence. Let $B=G\square K_2$, fix
$s\in V(B)$, and, for $A\subseteq V(B)$ and $x\in A$, let $F_s(A,x)$
be the number of self-avoiding walks from $s$ to $x$ whose vertex set
is exactly $A$. Set
\[
  F_s(\{s\},s)=1,
  \qquad
  F_s(A,s)=0\quad\text{for }A\ne\{s\},
\]
and set $F_s(A,x)=0$ if $s\notin A$. For $x\ne s$,
\begin{equation}
  F_s(A,x)=
  \sum_{\substack{y\in A\setminus\{x\}\\xy\in E(B)}}
  F_s(A\setminus\{x\},y).
  \label{eq:recurrence}
\end{equation}
Indeed, deleting the final vertex $x$ gives a bijection according to
the unique penultimate vertex $y$. Hence, if $N_k(s,t)$ denotes the
number of $k$-edge self-avoiding walks from $s$ to $t$, then
\begin{equation}
  N_k(s,t)=
  \sum_{\substack{A\subseteq V(B)\\s,t\in A,\ |A|=k+1}}
  F_s(A,t).
  \label{eq:length}
\end{equation}
Exact evaluation of \eqref{eq:recurrence}--\eqref{eq:length} gives
\[
\begin{array}{c|rrrrrrrrr|r}
 k&1&2&3&4&5&6&7&8&9&\mathrm{total}\\
\hline
 N_k(a_0,d_0)&1&0&2&1&9&14&13&13&7&60\\
 N_k(a_0,d_1)&0&2&0&7&9&12&16&9&4&59
\end{array}
\]
and therefore gives the claimed strict reverse inequality. All
computations use integer arithmetic.

It remains to verify minimum order. For an adjacent non-cut edge
$uv$, define
\[
  \Delta_G(u,v)=|\SAW(u_0,v_0)|-|\SAW(u_0,v_1)|.
\]
Up to isomorphism, the connected simple graphs on two, three, or four
vertices are exactly those listed below. The final column is the
maximum of $\Delta_G(u,v)$ over adjacent non-cut edges, evaluated by
the same recurrence; ``none'' means that every edge is a cut edge.
\[
\begin{array}{c|l|r}
\toprule
|V(G)|&G&\max \Delta_G(u,v)\\
\midrule
2&K_2&\text{none}\\
3&P_3&\text{none}\\
3&K_3&-1\\
4&P_4,\ K_{1,3}&\text{none}\\
4&C_4&-1\\
4&\text{paw graph}&0\\
4&K_4-e&-4\\
4&K_4&-14\\
\bottomrule
\end{array}
\]
The one-vertex case is vacuous. Thus no graph of order at most four
violates \eqref{eq:tang} at an adjacent non-cut edge.
\end{proof}

\section{A \texorpdfstring{$3$}{3}-connected counterexample}

The failure is not removed by imposing $3$-connectivity.

\begin{proposition}
Let $H$ be the graph in which $\{0,1,2,6\}$ induces $K_4$,
$\{3,4,5\}$ induces $K_3$, and the remaining edges are $05$, $14$,
and $23$. Then $H$ is $3$-connected and
\[
  |\SAW(0_0,5_0)|=5031>4801=|\SAW(0_0,5_1)|.
\]
\end{proposition}

\begin{proof}
Deleting at most two vertices leaves each clique nonempty and
connected. The three edges $05$, $14$, and $23$ are vertex-disjoint,
so at least one survives and joins the two remaining cliques. Hence
$H$ remains connected after deleting any two vertices, and therefore
$\kappa(H)\ge 3$. Since $\delta(H)=3$, we have $\kappa(H)=3$.
The exact counts follow from
\eqref{eq:recurrence}--\eqref{eq:length}.
\end{proof}

\begin{remark}
The examples above do not address Tang's complete-graph conjecture
\cite[Conjecture~1.8]{Tang}.
\end{remark}

\begin{thebibliography}{1}

\bibitem{Tang}
P.~Tang,
\emph{Maximum flow and self-avoiding walk on bunkbed graphs},
arXiv:2502.06237v1 [math.PR], 2025.

\end{thebibliography}

\end{document}